\section{Conclusion}
\label{sec:conclusion}

\methodname uses vector polynomial packing to reduce the dimension of
small-field relations before proving them over a large challenge
field. For inner and Hadamard products, the prover computes
row-polynomial products over \(\Field\) and reduces the original
length-\(N\) claim to a length-\(N/\kappa\) inner product over
\(\ExtensionField\).
Under the assumptions of Section~\ref{sec:highLevelEfficiency},
this gives \(O(\kappa N)\) base-field operations and \(O(N)\)
base-field coordinates of prover memory for the relation
reductions, excluding commitments and openings.
For \(|\Field|<\kappa\) and power-of-two \(\kappa\),
intermediate-field tabulation attains the same online Hadamard
bound with additional reusable tables. Section~\ref{sec:hadamard}
gives their preprocessing and storage costs.

The packed opening protocol authenticates challenge-dependent
evaluations against a fixed-basis witness commitment.
Together with the position-aware permutation argument, it yields
protocols for a single circuit execution and for multiple
executions of the same circuit.
The construction extends to binary arithmetizations by using
\(\mathbb F_2\subset\mathbb F_{2^8}\subset\ExtensionField\),
with byte-field arithmetic in the working-field PIOP.
Our component benchmarks show that the arithmetic and memory
savings are largest for the smallest witness fields tested.

Further work includes evaluating concrete PCS instantiations in
end-to-end proofs of quantized inference and adding compatible
zero-knowledge masking to the subprotocol transcripts.
The underlying reduction is modular: vector polynomial packing exploits
the witness field, and the terminal proof and commitment scheme
authenticate the resulting shorter relation.
